% methods.tex (v6): Methods with unnumbered subsections (Sci Rep). Numbers from numbers.tex.
\section{Methods}\label{sec:methods}

\bmhead{Study regions and observations}
Three regions with dense, multi-year thaw-depth observations were used (Fig.~1a): Alaska (\nncellAk{} cells, \nncyAk{} cell-years, \nnblkAk{} blocks of 0.5$^\circ$), the Lena Delta (\nncellLe{} cells, \nncyLe{} cell-years, \nnblkLe{} blocks) and north-western Canada (\nncellCa{} cells, \nncyCa{} cell-years, \nnblkCa{} blocks). Alaska observations are the ABoVE probe and ground-penetrating-radar thaw-depth collections\cite{Streletskiy2025,Talucci2024,Whitcomb2024}, restricted to August--October measurements with 0 < ALT < 300 cm. Lena Delta observations are the Samoylov and Kurungnakh thaw-depth series\cite{Boike2024}, restricted to August--September and to quality flags other than 2; for each location and year the seasonal maximum was taken. Canada and the remaining sites come from the Circumpolar Active Layer Monitoring (CALM) end-of-season grids\cite{Streletskiy2025}. Observations within the same 1 km cell and the same year were averaged into one cell-year label with its mean day of year; CALM end-of-season values without a date were assigned day 244 (31 August). Only years 2010--2024, for which yearly forcing exists, were kept. \nnmultiAk{}, \nnmultiLe{} and \nnmultiCa{} cells in the three regions have observations in four or more years.

\bmhead{Yearly forcing}
Monthly ERA5-Land 2 m temperature, snow depth and soil temperature (layer 1) for 2010--2024\cite{munozsabater2021_era5_land} were sampled at each cell with a land fallback of up to $\pm$5 grid cells. From these we computed, for every cell and year, the thawing degree-day sum (TDD), the freezing degree-day sum of the preceding winter (October--April), the hydrological-year snow-water equivalent, soil-temperature degree-days and monthly cumulative TDD. Three forcing variants were compared on identical cell-years: \emph{observation-day forcing}, in which TDD is accumulated to the observation day by linear interpolation of the monthly cumulative sum (the main variant); \emph{year forcing}, the full calendar-year sum; and \emph{climatological forcing}, the 2011--2024 mean of the yearly values, which mimics the inputs of existing products. The ratio of observation-day to full-year TDD has a median of 0.73 in Alaska, 0.81 in the Lena Delta and 0.89 in Canada, which quantifies how much of the thaw season the full-year sum overstates.

\bmhead{Static covariates and products}
Static covariates were terrain (elevation, slope, aspect components, topographic position and roughness from Copernicus DEM), soil (SoilGrids 2.0 organic carbon, clay, sand, bulk density and depth to bedrock\cite{Poggio2021}, NCSCD organic-carbon content and gelisol fraction\cite{Palmtag2022}), latitude and permafrost extent\cite{Westermann2024}. The main feature set contains eight climate variables plus terrain and soil (\emph{climate+terrain+soil}); a ladder of nested sets (climate only; +terrain; +soil; +MODIS vegetation and snow-off date; +high-resolution soil, vegetation and land-cover classes; all) was used to test which information carries the gain. Six published products were extracted at the label cells with the value of the label year where the product is yearly (ESA CCI v4 and v5\cite{Westermann2024}, Wei et al.\cite{Wei2026}, Liu et al.\cite{Liu2024}, Yi and Kimball\cite{YiKimball2019}) or the single published value otherwise (Aalto et al.\cite{Aalto2018}). Products were used only as comparison targets, never as inputs.

\bmhead{Physical models}
The Stefan model writes $\mathrm{ALT} = E\sqrt{\mathrm{TDD}}$ with one coefficient $E$\cite{Stefan1891}. \emph{Source-coefficient Stefan} (P0) takes $E$ fitted on the other two regions; \emph{recalibrated Stefan} (P1) fits $E$ by least squares on the training fold; an affine variant and the Kudryavtsev model with a fitted scale\cite{Kudryavtsev1974}, an edaphic Stefan with soil-dependent coefficients and a non-negative least-squares ensemble of the physical models were run as sensitivity cases (Supplementary Information).

\bmhead{Learned and combined models}
All learned models used CatBoost gradient boosting\cite{Prokhorenkova2018} (600 iterations, learning rate 0.03, depth 6, three seeds averaged) on the same features. \emph{Direct learning} (D0) predicts ALT from the covariates and forcing. \emph{Physics-input learning} (F1) adds the recalibrated Stefan prediction as a feature. \emph{Physics pseudo-labels} (D1) augments the training set with an equal number of unobserved cells, outside the scoring fold, labelled with the training-fold Stefan prediction. \emph{Physics plus residual learning} (R1) keeps the recalibrated Stefan prediction as an anchor and adds $\lambda$ times a CatBoost model of the residual $y - \hat{y}_P$; the main setting is $\lambda = 0.5$, with $\lambda \in \{0.25, 1\}$ and $\lambda$ chosen by inner block cross-validation as sensitivity cases. Random forest, a multilayer perceptron, TabM and FT-Transformer\cite{Gorishniy2021,Gorishniy2025,Holzmuller2024} were substituted for CatBoost in D0 and R1 to test the learner axis.

\bmhead{Validation splits}
Five splits were applied to each region\cite{Roberts2017,Ploton2020}: random five-fold by cell; spatial-block five-fold with 0.5$^\circ$ blocks assigned to balanced folds; year hold-out (each observation year scored by models trained on the other years); block-by-year, in which both the block and the year of the scored cell-years are unseen; and region hold-out, in which a region is scored by models trained on the other two regions with a 100 km buffer. Random split is reported as the optimistic reference, spatial block as the main result. Block sizes of 0.25$^\circ$ and 1$^\circ$ and alternative fold seeds were run as sensitivity cases.

\bmhead{Scoring}
Root-mean-square error (RMSE), bias and correlation were computed on out-of-fold predictions, together with a block-equal-weight RMSE and anomaly metrics after removing each cell's mean. Differences between two methods were computed on the same cell-years and their 95\% confidence intervals from 2000 block-bootstrap resamples of the 0.5$^\circ$ blocks. Gains were examined within terciles of regional TDD. Year-wise metrics were computed per observation year, and for cells observed in several years the correlation between observed and predicted year-to-year changes was computed.

\bmhead{Comparison with products}
Products were scored on the same cell-years as our out-of-fold predictions. Two conditions were applied: common cells, in which only cell-years with a value in every compared product are used, and a leakage mask, in which cells within 25 km (5 km as a sensitivity case) of the training sites of Wei et al. and Liu et al. are removed, because both products were trained on Alaska observations that overlap ours.

\bmhead{Yearly maps and area of applicability}
For Alaska and the Lena Delta the recalibrated Stefan and the physics-plus-residual models were fitted on all cell-years of the region and applied to a 1 km grid with the yearly ERA5-Land forcing of 2011--2024, giving one map per year and the 2011--2024 mean. The area of applicability was defined following Meyer and Pebesma\cite{MeyerPebesma2022}: a grid cell is inside when its nearest-neighbour distance to the training rows in standardised feature space does not exceed the 95th percentile of the distances between training rows of different blocks. Prediction intervals were taken from the 5th and 95th percentiles of block-cross-validation residuals.

\bmhead{Code availability}
All processing, modelling, evaluation and figure code is available at [DECISION: repository]. Experiments were run with fixed seeds; the yearly forcing, label and prediction tables needed to reproduce every number in the text are deposited with the data.
